\section{Algebry Hilberta}

Rozważamy unitarne reprezentacje grup zwartych $G$ z unormowaną miarą Haara $dx$, to znaczy $dx(G)=1$.

\begin{definition}{H-*-algebra}{}
  Algebra $A$ z inwolucją $x\mapsto x^*$ która jest przestrzenią Hilberta z iloczynem skalarnym $\langle \cdot,\cdot\rangle$ nazywa się \buff{H-*-algebrą} jeśli 
  \begin{itemize}
    \item $(\forall\;x,y,z\in A)\;\langle xy, z\rangle=\langle y, x^*z\rangle$
    \item $\|x^*\|=\|x\|$
  \end{itemize}
  Dodatkowo chcemy, żeby ta algebra była \acc{półprosta}, czyli dla $x\neq 0$ wymagamy $xx^*\neq 0$.
    %
    % Powiemy, że przestrzeń Hilberta $A$ z unitarnym działaniem $*$ nazywamy *-algebrą jeśli
    % \begin{itemize}
    %   \item $(\forall\;x,y,z\in A)\;\langle xy, z\rangle=\langle y, x^*z\rangle$
    %   \item $\|x^*\|=\|x\|$
    % \end{itemize}
\end{definition}

\begin{definition}{algebra Hilberta}{}
  W literaturze algebrą Hilberta nazywamy algebrę $A$ z inwolucją $x\mapsto x^*$ oraz iloczynem skalarnym $\langle \cdot,\cdot\rangle$ taką, że
  \begin{itemize}
    \item $(\forall\;x,y,z\in A)\;\langle xy, z\rangle=\langle y, x^*z\rangle$
    \item $(\forall\;x,y\in A)\;\langle x,y\rangle=\langle y^*,x^*\rangle$
    \item dla dowolnego ustalonego $a\in A$ operator $x\mapsto ax$ jest ograniczony
    \item przestrzeń liniowa rozpięta przez $xy$ dla $x,y\in A$ jest gęstym podzbiorem $A$.
  \end{itemize}
\end{definition}



